\section{Continuation-local synthesis}
\label{sec:continuations}

\subsection{A concrete composition failure}

Leant's current indexing diagnostics are unusually useful because they separate
missing expressiveness from poor composition order. The original bounded Haskell
query produces 256 checked candidates and falsifies all of them. Separate probes
find both function-valued Church folds and the required integer-case step; the
step appears at candidate 5 in an isolated context and candidate 219 in a larger
context. The report explicitly warns that neither number is the step's ordinal
inside the original whole-term search \cite{leant-diagnostic}.

The appropriate response is therefore not ``add function types to the carrier
grammar.'' They are already supported. It is to expose a fold's algebra as a
separate synthesis problem, constrain that algebra using the residual
specification, and avoid repeatedly discovering its local pieces inside many
unrelated whole terms.

\subsection{Recover the carrier from the elimination spine}

Suppose a candidate expression is obtained by eliminating a polymorphic value
and then applying the result to arguments of types $B_1,\ldots,B_k$, with final
result type $C$. One useful carrier candidate is
\begin{equation}\label{eq:carrier}
 R=B_1\to\cdots\to B_k\to C.
\end{equation}
This is a constraint implied by that particular application spine, not a general
solution to higher-rank type inference. Other spines, other carriers, and
alternative polymorphic instantiations remain separate search alternatives.

For integer indexing with a supplied default $d:A$, the last application is to
an index $i:\Z$, and the observed result has type $A$. The carrier suggested by
\eqref{eq:carrier} is $R=\Z\to A$. A right fold therefore needs
\[
 b:\Z\to A,
 \qquad a:A\to(\Z\to A)\to(\Z\to A).
\]
The nonstructural input has moved \emph{inside} the fold carrier. This is the
functional counterpart of generalizing an accumulator before induction, but the
search object is now a local algebra rather than a strengthened theorem.

\subsection{Derive contracts from a specification, not from a type}

Let $J_d(xs,i)$ denote the total indexing operation that returns $d$ at negative
or out-of-range indices. Its structural equations are
\begin{align}
 J_d([],i)&=d,\label{eq:index-base}\\
 J_d(h::xs,i)&=
 \begin{cases}
 d,&i<0,\\
 h,&i=0,\\
 J_d(xs,i-1),&i>0.
 \end{cases}\label{eq:index-step}
\end{align}
The type alone does not select this behavior: always returning $d$ has the same
result type. The equations must come from the user's specification, a proved
library characterization, or a separately proved unfolding of the intended
operation.

For the selected carrier, the equations induce local algebra contracts:
\begin{align}
 b(i)&=d,\label{eq:base-contract}\\
 a(h,k)(i)&=
 \begin{cases}
 d,&i<0,\\
 h,&i=0,\\
 k(i-1),&i>0,
 \end{cases}\label{eq:step-contract}
\end{align}
where $k$ is an arbitrary continuation. Treating $k$ as arbitrary is decisive.
A candidate step is not allowed to exploit the behavior of one guessed tail
function that happens to occur in a few examples.

\begin{theorem}[Fold contract]\label{thm:fold}
If $b$ and $a$ satisfy \eqref{eq:base-contract} and
\eqref{eq:step-contract}, then for every finite list $xs$ and integer $i$,
\[
 \operatorname{foldr}(a,b,xs)(i)=J_d(xs,i).
\]
\end{theorem}
\begin{proof}
For the empty list, the left side is $b(i)$, so the base contract proves the
claim. Suppose the claim holds for $xs$ at every integer index. The fold of
$h::xs$ is $a(h,\operatorname{foldr}(a,b,xs))$. For a negative index and for zero,
the step contract gives the first two branches of \eqref{eq:index-step}.
For a positive index it gives
$\operatorname{foldr}(a,b,xs)(i-1)$, which is $J_d(xs,i-1)$ by the induction
hypothesis. These cases exhaust the integers. The quantification over every
index in the induction hypothesis is exactly why the carrier is $\Z\to A$.
\end{proof}

The synthesized term has the following shape:
\begin{lstlisting}
fun default outerIndex input =>
  input (Int -> A)
    (fun head tailFn localIndex =>
      if localIndex < 0 then default
      else if localIndex = 0 then head
      else tailFn (localIndex - 1))
    (fun _ => default)
    outerIndex
\end{lstlisting}
This is explanatory Lean-like syntax. The archive's executable version uses a
small explicitly typed intermediate language, and its Lean specimen states the
corresponding theorem for encoded lists.

\subsection{The Church representation boundary}

For a finite list $xs$, define its Church encoding by
\[
 \operatorname{encode}(xs)(R,a,b)=\operatorname{foldr}(a,b,xs).
\]
Theorem~\ref{thm:fold} immediately proves the indexing theorem for
$\operatorname{encode}(xs)$. It does \emph{not} establish that every arbitrary
Lean value with a Church-style type behaves as the encoding of a finite list.
That stronger statement requires an appropriate representation or naturality
hypothesis and a proof in the actual source theory.

This matters especially in a synthesis engine that treats polymorphic values as
opaque data. A source type is not an automatic relational-parametricity theorem.
The prototype's acceptance language is deliberately named
\code{encoded-finite-list-index-v1}. Changing it to an unrestricted
Church-value claim causes the independent checker to reject the certificate.
The article's theorem and the executable checker therefore have the same scope.

\subsection{The executable search}

The prototype builds the fold sketch, but does not insert the answer's four
expressions as already chosen leaves. It enumerates four finite hole grammars:
the base result, the negative branch, the zero branch, and the positive branch.
The context contains the default, the input fold, the outer index, and $n$
unhelpful element-valued binders. The step additionally contains the head,
continuation, and local index. Candidate continuation arguments include
$i+2,i+1,i,i-2,i-3,i-1$ in that order.

The flat policy enumerates tuples from the four grammars, type-checks each whole
term, and evaluates the same six training examples. The local policy first
filters each hole against its local contract, using distinct symbolic-looking
values for element binders and a continuation observation $k(j)=({\tt tail},j)$;
it then enumerates the surviving tuples. Both policies use the same grammar
and candidate order within each hole.

Those local observations remain a search filter. Final acceptance uses an
independent symbolic checker: it checks the carrier, the base and step binders,
the outer application, the three branch results, and the affine predecessor
argument. In particular, the positive branch must call the locally bound
continuation at precisely $i-1$, not at the outer index. This checker is narrow
by design, but its success is stronger than passing a finite example suite.
The proof that the checked algebra implies list correctness is
Theorem~\ref{thm:fold}; its Lean formalization is a remaining integration gate.

\begin{lstlisting}
completeFold(demand, specification):
  for spine in compatibleEliminationSpines(demand):
    R := carrierEquation(spine)
    sketch := instantiateFoldAt(R)
    contracts := deriveBaseAndStepContracts(specification, sketch)
    for compatibleTuple in join(localSearch(contracts)):
      term := assemble(sketch, compatibleTuple)
      if sourceTypeCheck(term) and proveAlgebraContracts(term):
        return applyCheckedFoldTheorem(term)
  return Unknown
\end{lstlisting}

The implemented fragment has one selected fold schema, one carrier spine, and a
fixed three-way integer case split. Automatic discovery of arbitrary recursive
schemas, branch predicates, and source-level contracts is proposed rather than
implemented. The value of the prototype is to isolate and measure local
composition once these structures are available.

\subsection{Do not replace relational compatibility by projections}

Independent filtering is exact only when the specification factors into
independent local conditions. Suppose two holes $u,v$ must satisfy $u+v=1$ and
their domains are $\{0,1\}$. The allowed relation is
\[
 C=\{(0,1),(1,0)\}.
\]
Projecting onto the two coordinates and taking their Cartesian product introduces
$(0,0)$ and $(1,1)$. Both projections are individually correct, but their product
is not the original relation.

A general implementation should therefore store residual constraints as finite
relations or symbolic relations indexed by \emph{shared hole identifiers}.
For finite relations, use natural joins: if $R_1$ mentions holes $(u,v)$ and
$R_2$ mentions $(v,w)$, index both by the full identity of $v$, join matching
rows, and retain witnesses for all constraints. For symbolic relations, keep
shared metavariables owned by the same synthesis region and discharge the
conjoined residual condition before accepting an assembled term.

The indexing fragment factors after the constructor contract is fixed: the
base, negative, zero, and positive results have separate conditions. Its
one-survivor-per-hole result must not be generalized into a claim that all
higher-order synthesis admits such a product decomposition.

\subsection{Source identity, instantiation, and scheduling}

The proposed Leant/Djex integration needs one new search object: an
\emph{algebra region}. It retains the exact fold occurrence, selected source
polytype, instantiated carrier, source universes, available lexical binders,
local instance evidence, hole dependency relation, and residual contract. These
are references to the existing typed-source inventory, not fresh spellings that
will later be guessed back into the original context.

Each hole receives only the source binders its contract permits. Removing a
binder from a search context is a completeness-affecting choice unless a
separation argument justifies the smaller interface. Accordingly, the first
implementation should use a smaller interface as a preferred search branch,
while preserving a budgeted branch with the original context. The finite
prototype instead fixes the interface explicitly as part of its problem.

Search continuations should be memoized by the algebra region and hole, not by
pretty-printed type alone. A useful memo entry contains the next candidate
cursor, already considered local terms, remaining shared budget, and the residual
compatibility relation. Sibling holes may reuse a function head even when an
ancestor-path rule forbids repeating that same head along a single application
path; that distinction already matters in Djex's construction policies
\cite{djex-rank}.

This is not an invitation to reset the old global limits. Local checks, joins,
whole-term construction, and verification must all be charged. The intended
advantage is fewer useless complete combinations, not more hidden computation.
A cvc5 SyGuS worker could later solve first-order arithmetic holes inside a
fixed carrier, but it should not be asked to recover Lean's complete dependent
or rank-N source semantics \cite{cvc5}. The initial implementation needs no SMT
solver.

\subsection{What the experiment isolates}

With 12 distractor binders, the flat search reaches its first accepted term after
50,940 whole-term candidates. The local search checks 61 components and assembles
one whole term. A 256-candidate cutoff truncates the flat run, whereas the local
run succeeds. These numbers measure the supplied finite grammar, whose ordering
puts distractors early for both policies. They are not a measured 50,940-fold
speedup and not evidence that the original Leant query has been fixed.

The useful experimental conclusion is narrower and stronger: an explicit
constructor contract can remove the combinatorial cross-product that causes a
composition search to repeatedly combine locally wrong pieces. The production
test must preserve Leant's original source signature, primitive inventory,
binder order, behavioral predicate, and limits; the easier isolated step is a
regression fixture, not a substitute acceptance target.
